\documentclass[10pt]{article}

\usepackage[margin=0.82in]{geometry}
\usepackage{amsmath,amssymb,array,booktabs,microtype,multicol}
\usepackage[hidelinks]{hyperref}
\setlength{\emergencystretch}{1em}

\newcommand{\E}{\mathbb E}
\newcommand{\R}{\mathbb R}
\newcommand{\KL}{\mathrm{KL}}
\newcommand{\diag}{\operatorname{diag}}
\newcommand{\Span}{\operatorname{span}}

\title{How Protein Models Organize Biological Information:\\
A Multi-Axis Framework for Evolution, Structure, Function, and Design}
\author{Review working draft}
\date{July 2026}

\begin{document}
\maketitle

\begin{abstract}
Protein models are commonly classified by architecture or benchmark performance.  We instead ask
how they acquire, represent, compress, and reuse biological information.  Two statistical routes
dominate.  Query-specific systems construct an MSA and estimate the current family's evolutionary
distribution; amortized systems compress regularities from many families into pretrained language-
model parameters.  Within either route, information moves among sequence, residue-position,
categorical-state, pair, coordinate, and functional representations through explicit tracks and
bottlenecks.  We use this multi-axis framework to reinterpret AlphaFold 1--3, MSA Transformer,
RoseTTAFold, ESM, ESMFold, and ESM3.  We then show that protein design is the inverse or joint form
of the same statistical problem: RFdiffusion and ProteinMPNN factor backbone from sequence,
MultiFlow couples discrete and continuous flows, and Protpardelle and Pallatom fuse sequence,
side-chain chemistry, and all-atom geometry.  The central conclusion is that modern systems
increasingly separate computational representations without necessarily identifying the biological
causes that produced them.  Attention, pair states, modality tracks, and all-atom coordinates become
biologically interpretable only under declared categorical gauges, matched nulls, counterfactual
interventions, negative design, and independent experiments.
\end{abstract}

\section{Introduction}

The central problem of protein modeling is not merely to map sequence to structure.  It is to decide
which biological evidence is available, how that evidence is represented, and which information is
discarded before a prediction or design is produced.  Similar outputs can arise from very different
statistical experiments.  A query MSA measures one protein family's current sequence variation;
a protein language model recalls regularities previously learned from a global mixture of families;
a structure-trained network adds a prior over geometries observed in deposited structures.

These sources become difficult to distinguish after neural processing.  An MSA contains ancestry
and sampling variation across rows, conservation and gaps across columns, and dependence among
positions and residue states.  Attention, outer products, pair tracks, triangle operations, and
recycling combine these signals with templates and learned structural priors.  A representation can
therefore predict contacts or coordinates without being identifiable as direct coevolution or a
physical interaction.

We organize the field around a single information map:
\begin{equation}
\boxed{
\begin{gathered}
\{X_f,\theta_{\mathrm{LM}},\text{templates},\text{chemical context}\}
\longrightarrow
\text{sequence/pair/coordinate states}\\
\longrightarrow
\text{structure and function}
\mathrel{\substack{\longrightarrow\\[-4pt]\longleftarrow}}
\text{conditional protein design}.
\end{gathered}}
\label{eq:unified-information-map}
\end{equation}
The forward direction covers evolutionary inference and structure prediction.  The reverse or joint
direction covers inverse folding, backbone generation, sequence--structure co-design, and all-atom
generation.  The same questions apply in both directions: where does information enter, which state
stores it, how do states communicate, where is an axis compressed, which objective assigns meaning,
and which biological factors remain mixed.

This perspective distinguishes computational factorization from biological disentanglement.
Separate tensor axes, tracks, token alphabets, or prediction heads make a model easier to control and
analyze.  They do not prove that phylogeny, structural constraint, function, and chemical mechanism
occupy independent latent variables.  Stronger claims require matched nulls, interventions that vary
one source while controlling others, and validation against external biological measurements.

The review proceeds in four steps.  We first define a multi-axis coordinate system and an operational
criterion for disentanglement.  We then analyze two routes from evolutionary data to structure:
query-specific inference through an explicit MSA and amortized inference through pretrained language
models.  Next, we reinterpret protein design as alternative factorizations of the same joint
distribution.  Finally, we develop an interventionist research program and a technical outlook for
multistate, chemically typed, experimentally calibrated design.

\section{A unified information-factorization framework}

\subsection{Three axes of an alignment}

For $N$ homologous sequences, $L$ aligned positions, and alphabet size $q$, the one-hot MSA can be
written as
\begin{equation}
X\in
\mathcal H_{\mathrm{seq}}
\otimes
\mathcal H_{\mathrm{pos}}
\otimes
\mathcal H_{\mathrm{cat}}.
\label{eq:msa-product-space}
\end{equation}
Each axis has a different interpretation.
\begin{itemize}
\setlength{\itemsep}{1pt}
\item The \emph{sequence axis} contains ancestry, subfamilies, taxonomic imbalance, reweighting,
and finite-depth sampling.
\item The \emph{position axis} contains residue-indexed operators such as covariance, contact
scores, pair representations, and geometric constraints.
\item The \emph{categorical axis} contains amino-acid or nucleotide states and the distinction
between one-body preferences and multistate interaction contrasts.
\end{itemize}
A sequence principal component, a dominant eigenvector of a position matrix, and a categorical
constant mode are different objects.  AlphaFold architectures mix information across these axes,
but that mixing does not make the axes mathematically interchangeable.

\subsection{Query-specific and amortized evolution}

Let $f$ index a protein family and let $P_f(x)$ denote its evolutionary sequence distribution.
A query MSA provides a finite family-specific sample
\begin{equation}
X_f=\{x^{(1)},\ldots,x^{(N_f)}\},
\qquad x^{(n)}\sim P_f.
\label{eq:family-msa-sample}
\end{equation}
By contrast, a protein language model is trained on a global mixture
\begin{equation}
P_{\mathrm{train}}(x)
=\sum_f\pi_fP_f(x),
\label{eq:global-evolutionary-mixture}
\end{equation}
with family, taxonomic, and database sampling weights absorbed into $\pi_f$.  The model parameters
therefore amortize regularities shared across families.  At inference, a single query sequence must
implicitly identify which parts of this global prior are relevant.

We call $X_f$ \emph{query-specific evolutionary information} and the fitted language-model
parameters \emph{amortized evolutionary information}.  The distinction is not absolute: AlphaFold
also contains a learned cross-family prior, while ESM is trained on sequences whose redundancy and
phylogeny reflect the underlying database.  It identifies which evidence is newly acquired for the
query and which evidence is already stored in model weights.

\subsection{Marginal and interaction geometry}

Let $p_i(a)$ be the empirical or predicted marginal distribution at position $i$, and let
$D_i=\diag(p_i)$.  The weighted zero-mean projector is
\begin{equation}
P_i=I-\mathbf1p_i^{\mathsf T}.
\label{eq:categorical-projector}
\end{equation}
For a fixed categorical pair table $T_{ij}(a,b)$, the bivariate contrast is
\begin{equation}
G_{ij}=P_iT_{ij}P_j^{\mathsf T},
\qquad
p_i^{\mathsf T}G_{ij}=0,\quad G_{ij}p_j=0.
\label{eq:categorical-gauge}
\end{equation}
Under the product reference $p_i\otimes p_j$, this is the canonical functional-ANOVA projection
that removes constant and one-body terms \cite{hoeffding1948}.  The weighted strength is
\begin{equation}
m_{ij}=\|D_i^{1/2}G_{ij}D_j^{1/2}\|_F.
\label{eq:weighted-pair-strength}
\end{equation}
Equation~\ref{eq:categorical-gauge} is useful for interpreting MSA-derived couplings or learned
categorical fields.  It does not apply directly to an arbitrary hidden pair vector $z_{ij}$ unless
a decoder specifies how that vector acts on residue states.

\subsection{Four stages of an evolutionary structure model}

The systems can be compared through four transformations:
\begin{equation}
\boxed{
\text{sequence search}
\longrightarrow
\text{evolutionary representation}
\longrightarrow
\text{position-pair geometry}
\longrightarrow
\text{3D structure realization}.}
\label{eq:four-stage-pipeline}
\end{equation}
The major architectural choice is whether the evolutionary representation is constructed from a
query MSA or retrieved from a pretrained language model, and where it is compressed into pair
geometry.

The same distinction becomes operational in protein design.  Prediction estimates structure or
function from observed sequence, whereas design samples a sequence and structure under desired
constraints.  A design pipeline can explicitly factor backbone generation, inverse folding,
side-chain packing, and functional conditioning, or attempt to learn their joint distribution.
These factorizations determine both what can be controlled and which inconsistencies can arise at
the interfaces between models.

\subsection{An operational definition of disentanglement}

Let $B=(B_{\mathrm{evo}},B_{\mathrm{phy}},B_{\mathrm{chem}},B_{\mathrm{geom}},
B_{\mathrm{func}},B_{\mathrm{nuis}})$ denote evolutionary constraint, phylogeny, local chemistry,
three-dimensional geometry, function, and nuisance variation.  A learned representation
$r=T_\theta(B)$ is not disentangled merely because its covariance has several eigenvectors.  For
source $B_k$, operational disentanglement requires a declared readout or subspace $D_k(r)$ such that
\begin{equation}
D_k\!\left(T_\theta(\operatorname{do}(B_k=b))\right)
\text{ responds to }b,
\qquad
D_k\!\left(T_\theta(\operatorname{do}(B_\ell=b'))\right)
\text{ is controlled for }\ell\ne k.
\label{eq:operational-disentanglement}
\end{equation}
Because most biological factors cannot be independently randomized, practical evidence comes from
null MSAs, sequence interventions, modality dropout, architectural ablations, and external labels.

We distinguish four weaker architectural mechanisms.  \emph{Axis factorization} assigns sequence,
position, category, or coordinates to different tensor axes.  \emph{Track factorization} maintains
separate neural states such as MSA, pair, and coordinate tracks.  \emph{Objective factorization}
uses separate prediction heads or losses.  \emph{Bottleneck projection} deliberately removes an
axis, for example by averaging MSA rows into a pair state.  These mechanisms make analysis possible,
but none guarantees that phylogeny, conservation, structure, and function are statistically
independent inside a hidden channel.

For every model below we therefore ask the same questions: what information is observed; in which
state space is it stored; how states communicate; where information is compressed; which loss
assigns biological meaning; and which factors remain unidentifiable.

At the highest level, the systems separate into two evidence regimes:
\begin{center}
\small
\begin{tabular}{>{\raggedright\arraybackslash}p{0.17\linewidth}
                >{\raggedright\arraybackslash}p{0.30\linewidth}
                >{\raggedright\arraybackslash}p{0.39\linewidth}}
\toprule
Route & Evidence available at inference & Principal bottleneck \\
\midrule
Query-specific & current family sample $X_f$ plus learned structural priors & when row, subfamily, and categorical information is compressed into position-pair geometry \\
Amortized & query sequence or multimodal prompt interpreted through pretrained parameters & whether a global prior retains rare family constraints and separates sequence, structure, and function \\
\bottomrule
\end{tabular}
\end{center}
The following sections use representative models to examine these bottlenecks; the model-specific
audit is deferred to the synthesis rather than used as the organizing principle.

\section{Route I: query-specific evolutionary inference}

This route observes a finite sample $X_f$ from the query family during inference.  Its central
design choice is the lifetime of the sequence axis:
\begin{equation}
X_f
\xrightarrow{\text{family inference}}
r_f
\xrightarrow{\text{sequence-to-pair projection}}
z_{ij}
\xrightarrow{\text{geometric realization}}
a.
\label{eq:query-specific-route}
\end{equation}
Classical statistics and AF1 construct $r_f$ before the neural trunk.  AF2 keeps alignment rows as
an evolving state.  MSA Transformer explicitly factorizes row and column communication.  AF3
compresses a shallow MSA module into the pair track, while RoseTTAFold adds an active coordinate
track.  These systems differ less in whether they use evolution than in when family evidence becomes
irreversibly indistinguishable from learned geometry.

\subsection{Before AlphaFold: from covariance to learned pair potentials}

Classical coevolution methods convert an MSA into one- and two-position statistics.  Mutual
information and covariance measure marginal association but contain indirect effects.  Sparse
inverse covariance and direct-coupling analysis introduce probabilistic models intended to recover
conditional dependence \cite{jones2012,morcos2011,ekeberg2013,cocco2018}.  Their outputs are often
reduced to a scalar $L\times L$ contact score.

This reduction creates two losses of information.  First, a scalar block norm mixes the magnitude
of a categorical field with the marginal distributions under which it was estimated.  Second,
phylogeny, conservation, and finite-depth sampling can create low-rank position structure even
without genuine pair dependence.  APC and related corrections improve contact ranking but do not
provide a universal identification theorem for the removed mode \cite{dunn2008,wang2024}.

The immediate predecessors of AF1 replaced binary contacts with learned distance distributions.
Deep residual networks consumed profiles and coevolutionary features and predicted a distogram,
which supplies more geometric information than a contact threshold.  AlphaFold's first generation
combined this learned pair potential with a differentiable structure-realization procedure.

\subsection{AlphaFold 1: engineered MSA compression and distance potentials}

\subsubsection{Input and state space}

AF1 first constructs MSAs using sequence database search, including HHblits and PSI-BLAST
\cite{senior2020}.  Its released feature specification includes one-position quantities such as
HHM profiles, reweighted profiles, deletion probabilities, and alignment counts.  Pair inputs
include gap covariation, a $q^2$-dimensional pseudolikelihood coupling tensor, and the Frobenius norm
of that tensor.  The released CASP13 feature list contains, among others,
\begin{equation}
\texttt{hmm\_profile}_i,\quad
\texttt{reweighted\_profile}_i,\quad
\texttt{gap\_matrix}_{ij},\quad
\texttt{pseudolikelihood}_{ijab}.
\end{equation}
Thus AF1 is strongly MSA-based, but the neural trunk does not maintain the raw sequence axis.  The
alignment has already been compressed into site summaries and pair statistics.

In the language of Eq.~\ref{eq:msa-product-space}, AF1 applies an engineered map
\begin{equation}
X
\longmapsto
\bigl\{\phi_i(X),\Psi_{ij}(X)\bigr\}
\longmapsto
F_{\theta}\bigl(\phi_i,\phi_j,\Psi_{ij}\bigr),
\label{eq:af1-compression}
\end{equation}
where $\phi_i$ contains profile-like information and $\Psi_{ij}$ contains coevolutionary pair
features.  Any sequence-axis correction must therefore occur during feature construction; the
convolutional network cannot revisit individual homologous sequences.

\subsubsection{Communication, bottleneck, and output}

One-dimensional residue features are broadcast along the two axes of an $L\times L$ grid and
concatenated with pair features.  A deep dilated two-dimensional residual network predicts a
distribution over inter-residue distances.  Separate networks predict a background distance
distribution and torsion-related quantities.  The output is a distogram
\begin{equation}
P_{ij}(d_k\mid X),
\qquad k=1,\ldots,K,
\label{eq:af1-distogram}
\end{equation}
rather than a binary contact score.

The model converts the predicted distance distributions into a potential of mean force relative
to a learned background distribution.  Coordinates are then optimized by gradient descent, with
torsion and stereochemical terms supplying additional constraints.  The neural network predicts a
geometric energy landscape; it does not directly output the final atomic coordinates.

\subsubsection{Information retained and information mixed}

AF1 performs early sequence-axis compression and then learns on the position axis.  Its
pseudolikelihood tensor retains categorical state information, but scalar pair features such as
Frobenius norms are not marginal-orthogonal by construction.  Profile, entropy, gap, and coupling
strength can therefore remain entangled in the neural input.

The advantage is computational and statistical simplicity: engineered summaries expose known
coevolutionary signals to a powerful 2D network.  The limitation is that the model cannot learn a
new representation of phylogeny or alignment-row heterogeneity after compression.  AF1 should
therefore be described as an \emph{MSA-feature model}, not a raw-MSA neural architecture.

\subsection{AlphaFold 2: joint inference on MSA and pair states}

\subsubsection{Input and state space}

AF2 is an architectural discontinuity rather than a deeper AF1.  It embeds sampled MSA rows into
\begin{equation}
m_{si}\in\R^{c_m},
\end{equation}
where $s$ indexes homologous sequences and $i$ indexes residue positions.  In parallel it maintains
a pair representation
\begin{equation}
z_{ij}\in\R^{c_z}.
\end{equation}
The initial pair representation is formed from target-sequence embeddings and relative position;
templates and extra MSA sequences provide additional updates \cite{jumper2021}.
This explicit neural treatment of alignment rows is shared at a broad level with MSA language
models, although the Evoformer is distinguished by its coupled geometric pair track
\cite{rao2021}.

The main trunk contains 48 Evoformer blocks.  Each block performs
\begin{align}
m&\xrightarrow{\text{row attention with pair bias}}
\xrightarrow{\text{column attention}}
\xrightarrow{\text{transition}}m',\\
z&\xrightarrow{\text{outer-product mean from }m'}
\xrightarrow{\text{triangle multiplication/attention}}
\xrightarrow{\text{transition}}z'.
\end{align}
The pair state biases row-wise attention in the MSA, while the MSA updates the pair state through
the outer-product mean.  This creates repeated bidirectional communication between the sequence and
position axes.

\subsubsection{Communication and bottleneck}

The key sequence-to-pair operation has the schematic form
\begin{equation}
\Delta z_{ij}
=W\left[
\frac{1}{N}
\sum_{s=1}^N
a(m_{si})\otimes b(m_{sj})
\right].
\label{eq:outer-product-mean}
\end{equation}
This resembles a learned second moment, but the vectors $m_{si}$ have already been transformed by
attention, pair bias, and previous Evoformer blocks.  It is therefore not a fixed covariance or
Potts coupling.  The operation learns which categorical and sequence features should contribute to
the position-pair state.

AF2's own ablation study is informative.  Replacing raw MSA rows by only the profile and pairwise
frequency tables degrades performance.  This shows that access to individual alignment sequences
provides useful information beyond fixed first- and second-order summaries.  It does not prove
that every useful component is higher-order coevolution; sequence weighting, phylogeny, deletion
patterns, and adaptive denoising can also contribute.

\subsubsection{Output and supervision}

Triangle multiplicative updates and triangle attention enforce consistency among edges
$(i,j),(i,k),(j,k)$.  They treat $z$ as a dense relational geometry rather than an independently
scored contact map.  After the Evoformer, the first MSA row is projected to a single representation
$s_i$, while $z_{ij}$ is retained.

The structure module maps $(s_i,z_{ij})$ directly to coordinates.  Invariant point attention uses
pair features to bias communication among residue frames.  Each residue is represented by a rigid
backbone frame plus predicted torsion angles, and the module is trained end to end using frame
aligned point error and auxiliary geometric losses.  Recycling feeds the predicted structure,
pair representation, and first-row MSA representation back into the network for repeated
refinement.

\subsubsection{Information retained and information mixed}

AF2 delays sequence-axis compression.  The learned MSA tensor remains available throughout the
main trunk, so the model can adaptively combine row modes and position modes.  At the same time,
the representation is nonlinear, input-dependent, and repeatedly conditioned on $z$.  Its
attention spectra cannot be interpreted as phylogeny or direct coevolution from eigenvalue order
alone.

From the categorical perspective, $z_{ij}$ is a hidden vector, not an explicit $q\times q$ pair
potential.  Auxiliary masked-MSA and distogram heads encourage evolutionary and geometric content,
but they do not impose the marginal-orthogonal gauge in Eq.~\ref{eq:categorical-gauge}.  A decoded
categorical field would still require a stated reference and projection before it could be called
a one-body-free interaction.

AF2 is best described as a \emph{joint MSA--pair inference system}.  MSA information is not merely
an input feature; it is an evolving latent state coupled to structural geometry.

\subsection{AlphaFold 3: MSA-conditioned all-atom diffusion}

\subsubsection{Input and molecular scope}

AF3 predicts complexes containing proteins, RNA, DNA, small molecules, ions, modified residues,
and covalent links \cite{abramson2024}.  Its tokenizer assigns one token to a standard amino acid
or nucleotide and generally one token to each heavy atom of other chemical components.  This
broader domain already prevents a purely MSA-centric architecture: ligands and ions have no
homologous-sequence alignment.

The data pipeline nevertheless performs genetic search for polymer chains.  Protein chains are
searched against protein sequence databases; RNA chains are searched against Rfam, RNACentral,
and nucleotide collections.  AF3 constructs MSAs with up to 16,384 rows and uses paired rows where
cross-chain genetic information is available, extending the paired-MSA strategy developed for
AlphaFold-Multimer \cite{evans2022}.  If no hits are found, the query itself forms a
one-row MSA.  Protein templates remain optional and are searched per chain; non-protein ligands do
not receive sequence templates in this sense.

\subsubsection{Early sequence-axis bottleneck}

AF3 embeds MSA rows into $m_{si}$ but processes them in only four homogeneous MSA-module blocks.
Each block includes
\begin{equation}
m
\xrightarrow{\text{outer-product mean}}z,
\qquad
m\xrightarrow{\text{pair-weighted averaging}}m',
\qquad
z\xrightarrow{\text{triangle operations}}z'.
\label{eq:af3-msa-module}
\end{equation}
Unlike AF2, AF3 does not directly combine information across MSA rows by column-wise attention.
Alignment rows communicate through the pair representation.  The design explicitly aims to place
as much polymer information as possible into $z_{ij}$ before the main trunk.

After four blocks, the MSA representation is discarded and only the updated pair representation
is passed forward.  Thus AF3 still uses raw MSA rows, but the explicit sequence axis is no longer a
persistent state of the core network.

\subsubsection{Persistent pair state}

The Pairformer contains 48 blocks operating on a single token representation $s_i$ and pair
representation $z_{ij}$.  It retains the triangle operations of AF2 and uses pair-biased attention
to update the single representation.  There is no column-wise MSA attention.  Moreover, within
the Pairformer the single representation does not update the pair representation; information
flow from MSA and templates has already been deposited into $z$ by earlier modules.

This yields a one-way conditioning pattern in the core trunk:
\begin{equation}
z_{ij}
\longrightarrow
\text{bias and condition token communication},
\qquad
s_i\not\longrightarrow z_{ij}
\quad\text{inside Pairformer}.
\label{eq:pairformer-flow}
\end{equation}
The pair representation acts as the persistent relational memory of sequence, template, chemical,
and recycling information.

\subsubsection{Output and supervision}

AF3 replaces AF2's invariant-point-attention structure module with a conditional diffusion model
over all heavy-atom coordinates.  During training, Gaussian noise is added to the complete
structure and a denoiser predicts the clean coordinates conditioned on $s_i$ and $z_{ij}$.  The
same output space handles proteins, nucleic acids, ligands, ions, and modified components without
hard-coding a protein-specific rigid-frame parameterization.

The diffusion head is generative: multiple coordinate samples can be drawn from the same trunk
conditioning.  Most computation remains in the conditioning trunk, while the denoising network is
comparatively cheaper.  Confidence heads predict pLDDT, aligned error, and pairwise distance
confidence from the trunk and sampled structure.

\subsubsection{Information retained and information mixed}

AF3 makes the pair representation the main information bottleneck.  On proteins and RNA, $z$ is
MSA-conditioned; on ligands and ions, it must be constructed from chemical identity, covalent
structure, neighboring polymers, templates, and learned priors.  The same position-pair space
therefore contains heterogeneous evidence sources.

This heterogeneity makes a universal ``coevolution mode'' interpretation untenable.  A pair mode
linking two protein residues may contain evolutionary information, while a protein--ligand pair
cannot.  Spectral analysis of AF3 must condition on entity type and input provenance.  The correct
description is an \emph{MSA-conditioned multimolecular pair model}, not a purely MSA-based model.

\subsection{MSA Transformer: separating alignment axes without discarding the MSA}

\subsubsection{Input, state space, and objective}

MSA Transformer is trained across many protein families, but each inference input is an explicit
alignment $X_f$ \cite{rao2021}.  After embedding, its hidden state remains
\begin{equation}
H\in\R^{N\times L\times d}.
\end{equation}
Each layer alternates column attention across homologous sequences at fixed position, row attention
across residue positions, and a channel-wise transition:
\begin{equation}
H\xrightarrow{\operatorname{Attn}_{\mathrm{seq}}}
H^{c}\xrightarrow{\operatorname{Attn}_{\mathrm{pos}}}
H^{r}\xrightarrow{\operatorname{MLP}}H'.
\label{eq:msa-transformer-axes}
\end{equation}
The row-attention map is tied across MSA rows, imposing the hypothesis that homologues share one
position--position structural organization.  The model minimizes masked reconstruction over MSA
entries,
\begin{equation}
\mathcal L_{\mathrm{MSA\text{-}MLM}}
=-\E_{X,M}\sum_{(s,i)\in M}
\log p_\theta(X_{si}\mid X_{\setminus M}).
\label{eq:msa-transformer-loss}
\end{equation}
Unlike a Potts model fitted separately to one family, the inference rule is shared across millions
of MSAs.  Unlike single-sequence ESM, the current family sample is still directly observed.

\subsubsection{Information retained and information mixed}

The architecture explicitly factorizes two computational axes.  Column attention can compare
residue states, gaps, and subfamily occupancy across sequences at one aligned position.  Tied row
attention forms a shared $L\times L$ routing operator that can accumulate covariation and structural
compatibility across homologues.  This is a genuine axis-level separation and explains why contact
patterns emerge in row attention without structure supervision.

It is not a biological disentanglement theorem.  Phylogenetic relatedness, redundancy,
conservation, gap patterns, and compensatory substitutions all enter the same $H_{si}$ channels and
are optimized by the same reconstruction loss.  Tying row attention assumes common geometry across
the family and can suppress subfamily-specific conformations.  Contact decoding from attention
still requires symmetrization, APC, and a logistic readout; it does not turn attention into a Potts
coupling.  MSA Transformer is therefore best described as a \emph{query-specific plus amortized
evolution model}: $X_f$ supplies current family evidence, while $\theta$ supplies a cross-family
denoising prior.

\subsection{RoseTTAFold: coupled sequence, pair, and coordinate tracks}

\subsubsection{Input, tracks, and communication}

RoseTTAFold organizes inference into three communicating representations \cite{baek2021}:
\begin{equation}
m_{si}\in\mathcal H_{\mathrm{seq}}\otimes\mathcal H_{\mathrm{pos}},
\qquad
z_{ij}\in\mathcal H_{\mathrm{pos}}^{\otimes2},
\qquad
x_i\in\R^3.
\label{eq:rosettafold-tracks}
\end{equation}
The 1D track processes the query MSA, the 2D track represents residue distances and orientations,
and the 3D track updates backbone coordinates with an SE(3)-equivariant transformer.  Information
is repeatedly exchanged rather than passed through only once:
\begin{equation}
m\rightleftarrows z\rightleftarrows x,
\qquad m\rightleftarrows x.
\label{eq:rosettafold-communication}
\end{equation}
Templates can enter as additional geometric evidence.  End-to-end coordinate and geometric losses
allow errors in the 3D track to modify how the 1D and 2D tracks extract MSA information.

\subsubsection{Information retained and information mixed}

RoseTTAFold makes the sequence/pair/coordinate distinction explicit.  This is stronger than AF1's
early feature compression because geometry remains an active state, and it differs from AF2's
Evoformer-to-structure ordering because approximate coordinates can communicate back during
three-track reasoning.  SE(3) equivariance removes dependence on the arbitrary global coordinate
frame, which is a rigorous disentanglement of a geometric nuisance symmetry.

The tracks do not isolate biological causes.  The 2D state combines MSA covariation, templates,
current coordinates, residue identity, and learned structural priors.  The 3D loss encourages every
predictive signal to cooperate, not to remain independently readable.  Consequently a pair mode
that predicts a contact cannot be assigned uniquely to direct coevolution, template evidence, or
geometric feedback without track ablation.  RoseTTAFold fits sequence compatibility, pairwise
distance/orientation geometry, and coordinate consistency; it does not explicitly fit function or
phylogeny as separate targets.

\section{Route II: amortized language-model inference}

This route moves the expensive evolutionary aggregation from query time to pretraining:
\begin{equation}
P_{\mathrm{train}}
\xrightarrow{\text{pretraining}}
\theta_{\mathrm{LM}},
\qquad
x\xrightarrow{\theta_{\mathrm{LM}}}
r_x\xrightarrow{\text{geometric or multimodal decoder}}(a,f).
\label{eq:amortized-route}
\end{equation}
The query-family sequence axis is never instantiated, but evolutionary information remains present
through $\theta_{\mathrm{LM}}$.  ESM stops at contextual sequence representations, ESMFold maps
those representations to coordinates, and ESM3 adds structure and function token tracks.  The gain
is rapid inference and transfer across families; the cost is weaker access to rare, current-family,
and subfamily-specific statistics.

\subsection{ESM: amortizing evolutionary information}

\subsubsection{Input, state space, and objective}

ESM models are trained on large collections of unaligned natural protein sequences rather than on
one MSA for each query \cite{rives2021}.  Let $x=(x_1,\ldots,x_L)$ and let $M$ be a random set of
masked positions.  The masked language-model objective is
\begin{equation}
\mathcal L_{\mathrm{MLM}}(\theta)
=-\E_{x\sim P_{\mathrm{train}}}
\E_{M\sim\nu}
\sum_{i\in M}
\log p_{\theta}(x_i\mid x_{[L]\setminus M}).
\label{eq:esm-mlm-objective}
\end{equation}
When one position is masked at a time, this resembles a nonlinear pseudolikelihood objective.  In
practice several positions are corrupted simultaneously, so the model learns denoising
conditionals under a distribution of missing contexts \cite{rao2020}.

The objective does not explicitly separate conservation, phylogeny, structure, and function.  Any
regularity that helps predict a missing amino acid can reduce the loss.  The learned conditional
may therefore contain
\begin{equation}
\log p_{\theta}(x_i\mid x_{-M})
=c_i+u_i(x_i)+
\sum_{j\notin M}g_{ij}(x_i,x_j)
+R_i^{(\geq2)}(x_i,x_{-M}),
\label{eq:esm-conditional-decomposition}
\end{equation}
where the displayed terms are a conceptual functional decomposition, not separate ESM parameter
blocks.  The Transformer can represent higher-order and context-dependent effects through its
depth and nonlinear attention.

Unlike a family-specific Potts model, ESM minimizes Eq.~\ref{eq:esm-mlm-objective} over the mixture
in Eq.~\ref{eq:global-evolutionary-mixture}.  Its conditionals are shared across protein families.
This parameter sharing explains why single-sequence language models can be particularly useful for
proteins with shallow MSAs, but it also means that rare family-specific constraints can be averaged
with broader sequence regularities.

\subsubsection{Communication and position operators}

For layer $\ell$ and head $h$, ESM forms a row-stochastic attention matrix
\begin{equation}
A_{ij}^{\ell h}
=\operatorname{softmax}_{j}
\left(
\frac{q_i^{\ell h\,\mathsf T}k_j^{\ell h}}
{\sqrt d}
\right).
\label{eq:esm-attention}
\end{equation}
Attention permits every residue to condition on every other residue.  Sparse linear combinations
of symmetrized attention heads, followed by APC in the published contact pipeline, recover
long-range contact information without a query MSA \cite{rao2020}.  ESM-1b can outperform a
family-specific pseudolikelihood pipeline on contact prediction, especially when the available MSA
is shallow.

This does not make $A^{\ell h}$ a direct biological interaction matrix.  It is normalized by
softmax, depends on the input sequence and layer context, and may route information needed for
motifs, syntax, position, or optimization rather than physical contact.  Its invariant row-sum
mode and non-normal dynamics also make eigenvalue deletion difficult to interpret.  Contact
readout is a supervised or lightly supervised decoder of attention, not an identity between
attention and contact.

ESM-MSA-1b is a separate branch that takes an explicit MSA and performs row and column attention
\cite{rao2021}.  It is useful as a bridge between the two paradigms, but the ESM-2 model used by
ESMFold is a single-sequence language model.

\subsubsection{Biological signals fitted and mixed}

The MLM loss rewards several nested forms of predictability.
\begin{enumerate}
\setlength{\itemsep}{1pt}
\item \textbf{Marginal composition and conservation.}  Amino-acid frequencies, low-complexity
regions, signal peptides, transmembrane composition, and family-specific conservation reduce token
uncertainty.
\item \textbf{Local sequence grammar.}  Motifs, secondary-structure propensities, turn patterns,
and local biochemical compatibility are visible in short contexts.
\item \textbf{Long-range structural constraints.}  Contacts, fold topology, disulfide pairing,
and compensatory substitutions can improve prediction across distant positions.
\item \textbf{Function and cellular context.}  Active-site motifs, binding families, localization
signals, and domain architecture correlate with sequence and can emerge in representations.
\item \textbf{Mutational tolerance.}  The log-odds score
\begin{equation}
\Delta_i(a\to b)
=\log p_{\theta}(b\mid x_{-i})
-\log p_{\theta}(a\mid x_{-i})
\label{eq:esm-mutation-score}
\end{equation}
is used as a zero-shot proxy for variant effect \cite{meier2021}.
\item \textbf{Phylogeny and database composition.}  Taxonomic redundancy and family abundance are
part of $P_{\mathrm{train}}$ unless explicitly reweighted.  ESM learns from these signals but does
not identify them as nuisance.
\end{enumerate}

Scaling improves the linear recoverability of secondary structure, contacts, function, and other
properties \cite{rives2021,lin2023}.  This is evidence that the representations contain these
signals, not that the MLM objective uniquely separates them.

\subsection{ESMFold: from amortized sequence statistics to coordinates}

\subsubsection{Input and state initialization}

ESMFold uses a pretrained ESM-2 model and requires only the query sequence at inference
\cite{lin2023}.  Let $h_i^{(\ell)}$ be the residue representation at ESM-2 layer $\ell$.  A learned
softmax combination initializes the folding single state,
\begin{equation}
s_i^{(0)}
=\operatorname{MLP}
\left(
\sum_{\ell=0}^{L_{\mathrm{LM}}}
\alpha_{\ell}h_i^{(\ell)}
\right)
+e(x_i),
\qquad
\alpha=\operatorname{softmax}(w).
\label{eq:esmfold-single-init}
\end{equation}
The attention maps from every language-model layer and head are concatenated and projected to form
the initial pair state,
\begin{equation}
z_{ij}^{(0)}
=\operatorname{MLP}
\left(
\operatorname{concat}_{\ell,h}
A_{ij}^{\ell h}
\right).
\label{eq:esmfold-pair-init}
\end{equation}
The released implementation freezes ESM-2 while training the folding trunk and output heads.  The
single state therefore receives amortized sequence statistics through hidden representations, and
the pair state receives a learned readout of position-to-position attention.

\subsubsection{Communication, bottleneck, and output}

The folding trunk contains 48 triangular self-attention blocks.  Each block updates the sequence
state using pair-biased attention, updates the pair state from the sequence state, and applies
triangle multiplication and triangle attention:
\begin{align}
s&\leftarrow s+\operatorname{Attention}(s;\operatorname{bias}(z))
+\operatorname{MLP}(s),\\
z&\leftarrow z+\operatorname{SeqToPair}(s)
+\operatorname{Triangle}(z)+\operatorname{MLP}(z).
\label{eq:esmfold-trunk}
\end{align}
Relative position and recycled distance bins provide additional pair information.  Up to four
recycling iterations repeatedly refine $(s,z)$ and the predicted structure.

The final states enter an AlphaFold 2-style structure module derived from OpenFold.  The module
predicts residue frames, torsions, atom coordinates, distograms, pTM, and pLDDT-related quantities.
Thus ESMFold replaces AF2's query MSA and Evoformer with ESM-2 plus a single-sequence folding trunk,
while retaining the dense pair geometry and structure-module logic.

\subsubsection{Sequence and structural supervision}

ESMFold combines two objectives trained on different data distributions.  First, ESM-2 fits the
sequence mixture through Eq.~\ref{eq:esm-mlm-objective}.  Second, the folding parameters $\phi$ fit
structure supervision,
\begin{equation}
\phi^*
=\arg\min_{\phi}
\E_{(x,y)\sim P_{\mathrm{struct}}}
\mathcal L_{\mathrm{fold}}
\left(
y,
F_{\phi}(s^{(0)}(x),z^{(0)}(x))
\right),
\label{eq:esmfold-structure-objective}
\end{equation}
where $y$ is an experimentally determined or distilled structure.  The language model attempts to
capture sequence regularities; the folding trunk learns which of those regularities predict
geometry.

Biologically, ESMFold attempts to reconstruct local stereochemistry, secondary structure,
long-range contacts, fold topology, and confidence from a single sequence plus the amortized prior
stored in ESM-2.  It does not observe the query family's current MSA, so it cannot directly measure
family-specific covariation, subfamily occupancy, or alignment depth.  Its speed comes from
replacing query-time evolutionary search with a shared pretrained inference map.

\subsubsection{Information retained and information mixed}

ESMFold's initial pair spectrum is inherited from the collection of language-model attention maps,
then transformed by 48 nonlinear pair blocks and recycling.  It should therefore be interpreted as
a learned structural compatibility geometry, not as a direct-coupling matrix.  A leading mode can
contain residue identity, secondary structure, relative position, fold class, confidence, and
training-set priors.

The central comparison is
\begin{equation}
\begin{aligned}
\text{AlphaFold:}\quad
&F_{\theta}(x,X_f,\text{templates}),\\
\text{ESMFold:}\quad
&G_{\phi}(x;\theta_{\mathrm{LM}}),
\end{aligned}
\label{eq:af-esmfold-comparison}
\end{equation}
where $X_f$ is a query-specific family sample and $\theta_{\mathrm{LM}}$ is an amortized global
evolutionary prior.  Neither route is intrinsically more ``biological''; they observe different
statistics and fail in different regimes.

\subsection{ESM3: multimodal sequence--structure--function generation}

\subsubsection{Input tracks, latent fusion, and objective}

ESM3 extends masked modeling from amino-acid sequences to several aligned token tracks
\cite{hayes2024}.  At residue $i$, the observed variables can include sequence $a_i$, a discrete
structure token $u_i$, secondary structure, solvent accessibility, function keywords, and
residue-level annotations.  Learned embeddings from the available tracks are added and processed
in one bidirectional Transformer latent state:
\begin{equation}
h_i^{(0)}
=\sum_{t\in\mathcal T}e_t(y_i^t)+e_{\mathrm{pos}}(i),
\qquad
h^{(L)}=\operatorname{Transformer}_\theta(h^{(0)};x^{\mathrm{3D}}).
\label{eq:esm3-track-fusion}
\end{equation}
When coordinates are supplied, invariant geometric attention in the first block conditions the
latent representation on local residue frames.  Structure tokens are produced by a discrete
autoencoder that compresses local atomic neighborhoods; a separate decoder reconstructs atomic
coordinates from generated structure tokens.

Training masks variable fractions of one or more tracks and predicts their missing tokens:
\begin{equation}
\mathcal L_{\mathrm{ESM3}}
=-\E_{y,M}\sum_{t\in\mathcal T}\sum_{i\in M_t}
\log p_\theta(y_i^t\mid y_{\setminus M}).
\label{eq:esm3-loss}
\end{equation}
This trains many conditional directions in one model: sequence from structure and function,
structure from sequence, function-compatible completion, and unconditional joint generation.

\subsubsection{Biological distribution fitted}

ESM3 attempts to fit the observational joint distribution
\begin{equation}
P_{\mathrm{train}}(
\text{sequence},\text{structure},\text{function},\text{annotations}).
\label{eq:esm3-joint-distribution}
\end{equation}
This distribution contains natural evolutionary selection, predicted structures, computational
function annotations, inverse-folding augmentations, and database sampling.  Conditional
generation can enforce a fold, catalytic annotation, secondary-structure pattern, or partial
sequence while completing the remaining variables.  It therefore targets sequence viability,
local and global geometry, biochemical function labels, solvent exposure, and compatibility among
these properties.

The phrase ``evolutionary simulation'' must be interpreted statistically.  ESM3 samples from a
learned distribution shaped by surviving natural proteins and augmented labels; it does not model
mutation, population size, selection coefficients, genetic drift, or historical time as explicit
dynamical variables.

\subsubsection{Information retained and information mixed}

ESM3 has the clearest objective-level factorization in this review: modalities have separate input
tokens, corruption masks, and output heads.  Modality dropout and cross-track prompting provide
direct tests of whether one track contains information predictive of another.  Structure
tokenization also creates an explicit bottleneck between atomic geometry and discrete local motifs.

Yet Eq.~\ref{eq:esm3-track-fusion} shows that the tracks are fused by addition into one latent
space.  A hidden direction need not be uniquely ``sequence,'' ``structure,'' or ``function.''
Moreover, structure and function labels are correlated with sequence family, and many are predicted
rather than experimentally measured.  Successful conditional generation establishes compatibility
under the training distribution, not causal independence.  ESM3 therefore provides
\emph{controllable multimodal factorization at the interface}, but not identifiable biological
disentanglement in its internal representation.

\section{Synthesis: computational factorization and biological entanglement}

The case studies reduce to two independent architectural choices.  The first is the source of
evolutionary evidence: a current family sample $X_f$, an amortized parameter prior, or both.  The
second is the representation bottleneck: when sequence rows, categorical states, pair geometry,
and coordinates cease to be separately accessible.  Performance can improve while identifiability
decreases because later states integrate more evidence but preserve less provenance.

\subsection{Where the sequence axis disappears}

The most informative comparison is the lifetime of $\mathcal H_{\mathrm{seq}}$:
\begin{align}
\text{AF1:}\quad
&X
\xrightarrow{\text{profiles/couplings}}
\Psi_{ij}
\xrightarrow{\text{2D ResNet}}
P_{ij}(d),\\
\text{AF2:}\quad
&X
\xrightarrow{\text{embedding}}
m_{si}
\mathrel{\overset{\text{outer-product mean}}
{\underset{\text{pair bias}}{\rightleftarrows}}}
z_{ij}
\xrightarrow{\text{structure module}}
x,\\
\text{AF3:}\quad
&X
\xrightarrow{\text{four-block MSA module}}
z_{ij}
\xrightarrow{\text{Pairformer}}
(s_i,z_{ij})
\xrightarrow{\text{diffusion}}
x,\\
\text{MSA Transformer:}\quad
&X\xrightarrow{\text{row/column attention}}H_{si}
\xrightarrow{\text{row-attention readout}}A_{ij},\\
\text{RoseTTAFold:}\quad
&(X,\text{templates})\xrightarrow{\text{three-track blocks}}
(m_{si},z_{ij},x_i),\\
\text{ESM:}\quad
&x
\xrightarrow{\theta_{\mathrm{LM}}}
\{h_i^{(\ell)},A_{ij}^{\ell h}\},\\
\text{ESMFold:}\quad
&x
\xrightarrow{\theta_{\mathrm{LM}}}
(s_i^{(0)},z_{ij}^{(0)})
\xrightarrow{\text{folding trunk}}
x^{\mathrm{3D}},\\
\text{ESM3:}\quad
&y_{\mathrm{observed}}^{\mathcal T}
\xrightarrow{\theta_{\mathrm{MM}}}h_i
\xrightarrow{\text{track heads/decoder}}
y_{\mathrm{missing}}^{\mathcal T},x^{\mathrm{3D}}.
\end{align}
AF1 compresses before the neural trunk, AF2 retains the axis throughout the trunk, and AF3
compresses it near the trunk entrance.  MSA Transformer retains the alignment tensor but
factorizes communication along its two axes.  RoseTTAFold retains an MSA track while adding active
pair and coordinate tracks.  ESM, ESMFold, and ESM3 never instantiate a query-family sequence axis
at inference.  Instead, population-level regularities were compressed during pretraining into
shared parameters.  Thus ``no MSA input'' means no query-specific empirical sample, not absence of
evolutionary information.

\subsection{What each system learns on the position axis}

AF1 predicts an explicit distribution over distances.  Its pair object has a direct geometric
readout.  AF2 and AF3 instead maintain hidden pair vectors with triangle operations.  These vectors
are more expressive but less directly interpretable.  Their spectrum reflects a learned metric of
structural compatibility, not a single contact or coupling matrix.

ESM creates position--position attention matrices $A^{\ell h}$ as routing operators for contextual
prediction.  Their symmetrized, APC-corrected combinations can predict contacts, but an attention
edge is not itself a contact or direct coupling.  ESMFold projects the full collection of these
maps into $z_{ij}^{(0)}$ and then applies triangle updates.  Its final pair geometry is therefore a
learned transformation of attention-derived compatibility, rather than a raw language-model
attention spectrum.

MSA Transformer obtains a position operator by tying row attention across homologues.  RoseTTAFold
makes position-pair geometry an explicit 2D track and exchanges it with coordinates.  ESM3 instead
compresses local atomic neighborhoods into categorical structure tokens; global geometry is only
recovered after contextual modeling and structure decoding.  These are respectively an attention
operator, a persistent geometric state, and a discrete geometric code, so their spectra are not
directly interchangeable.

The triangle operations can be viewed as nonlinear message passing on triples of positions.  They
enforce relational consistency but should not be confused with graph-Laplacian diffusion.  The
updates are channel-dependent, gated, and input-dependent; their operators are generally
non-normal, so singular values and Jacobian modes may be more informative than eigenvalues of one
attention matrix.

\subsection{What happens to categorical state geometry}

AF1 directly receives categorical profiles and Potts-like pair couplings.  AF2 embeds amino-acid
states into MSA channels and only exposes categorical predictions through auxiliary heads.  AF3
extends token categories beyond biological alphabets to arbitrary chemical components.  Across
the generations, categorical information becomes more latent and more heterogeneous.

ESM is trained by explicit categorical prediction: at masked positions it estimates a conditional
distribution over amino acids.  Nevertheless, its attention matrices live only on the position
axis, and its hidden channels are not an identifiable $q\times q$ interaction table.  ESMFold
inherits these categorical conditionals only indirectly through frozen ESM-2 representations;
its pair state is optimized for geometry, not for recovering a Potts gauge.

MSA Transformer predicts categorical residues at every MSA entry while sharing position maps across
rows.  RoseTTAFold embeds residue categories in its MSA track but supervises pair and coordinate
tracks geometrically.  ESM3 has explicit vocabularies for each modality, including amino acids,
structure codes, and function annotations; however, their embeddings are summed into a common
latent state.  Separate output alphabets therefore provide objective factorization without latent
independence.

For mechanistic interpretation, a hidden pair state should be decoded into a typed field
\begin{equation}
T_{ij}(a,b)=D_{\theta}(z_{ij};a,b)
\end{equation}
and projected using Eq.~\ref{eq:categorical-gauge}.  Without this step, a scalar norm or leading
mode of $z_{ij}$ can mix conservation, residue identity, chain type, geometry, confidence, and
interaction information.

\subsection{Disentanglement audit by model}

The architectures can now be compared without treating a named track as proof of biological
identifiability.

\begin{center}
\scriptsize
\begin{tabular}{>{\raggedright\arraybackslash}p{0.09\linewidth}
                >{\raggedright\arraybackslash}p{0.26\linewidth}
                >{\raggedright\arraybackslash}p{0.25\linewidth}
                >{\raggedright\arraybackslash}p{0.29\linewidth}}
\toprule
Model & Explicit separation & Main compression or bottleneck & Information that remains entangled \\
\midrule
AF1 & engineered one-site, pair, gap, and template features & raw MSA rows removed before the neural trunk & phylogeny, conservation, coupling magnitude, and finite-depth noise in scalar pair features \\
AF2 & persistent MSA and pair tracks; auxiliary categorical and geometric heads & outer-product mean projects sequence rows into $z_{ij}$ & MSA ancestry, residue identity, templates, pair geometry, recycling, and structural prior \\
AF3 & short MSA module, token and pair tracks, entity-aware all-atom output & MSA axis discarded before Pairformer; all evidence deposited in $z_{ij}$ & evolutionary, template, chemical, interface, confidence, and diffusion-conditioning signals \\
MSA Transformer & row and column attention factorize MSA axes & tied row map pools homologues into one position operator & phylogeny, redundancy, gaps, conservation, covariation, and family heterogeneity \\
RoseTTA\hspace{0pt}Fold & 1D MSA, 2D pair, and 3D coordinate tracks; SE(3) symmetry handling & repeated cross-track projections & coevolution, templates, coordinate feedback, and learned fold priors inside the pair track \\
ESM & categorical MLM head and layer/head organization & global families compressed into $\theta_{\mathrm{LM}}$; query compressed into hidden states & family identity, phylogeny, structure, function, motifs, and database composition \\
ESMFold & frozen LM, single and pair states, folding trunk, geometric heads & attention maps projected into $z_{ij}^{(0)}$ and repeatedly recycled & LM priors, residue identity, relative position, geometry, fold class, and confidence \\
ESM3 & separate modality tokens, masks, heads, and structure tokenizer & track embeddings summed into one latent state; geometry discretized & sequence family, predicted structure, function labels, annotations, and dataset augmentation \\
\bottomrule
\end{tabular}
\end{center}

The strongest exact separation in the table concerns representation type or symmetry: MSA rows
versus residue positions, pair states versus coordinates, or invariance to rigid motion.  The
separation of evolutionary history, physical constraint, and biological function is weaker because
these variables are correlated in natural data and are usually optimized jointly.

\subsection{MSA depth and amortized structural priors}

All three AlphaFold generations, MSA Transformer, and RoseTTAFold benefit from deeper informative
alignments, but none is merely a deterministic coevolution estimator.  Training across many
families or structures supplies a strong
prior over local geometry, folds, interfaces, and chemical environments.  Templates provide an
additional route.  AF2 and AF3 can therefore make useful predictions for shallow or query-only
MSAs, although uncertainty and failure rates generally increase.  MSA-free structure models such
as ESMFold make the complementary point that a sufficiently trained protein language model can
carry a substantial evolutionary prior in its parameters \cite{lin2023}.

The relationship can be summarized as
\begin{equation}
\begin{aligned}
\widehat x_{\mathrm{AF}}
&=F_{\theta}(x,X_f,\text{templates},\text{chemical graph},\text{recycling}),\\
\widehat x_{\mathrm{ESMFold}}
&=G_{\phi}\!\left(x;\theta_{\mathrm{LM}},\text{recycling}\right),
\end{aligned}
\end{equation}
where the learned parameter $\theta$ carries information accumulated across the structural
training distribution, while $\theta_{\mathrm{LM}}$ carries regularities accumulated across the
global sequence distribution.  Removing MSA input does not remove evolutionary information that
is already encoded statistically in trained parameters.

\section{Protein design as inverse and joint inference}

Prediction models approximate conditionals such as $p(a,f\mid x,X_f,\theta)$; design requires the
reverse conditional or a joint model under a specification $c$.  The apparent diversity of protein
design algorithms therefore reflects different choices for factorizing one joint distribution, not
a separate conceptual field.  Structure-first pipelines expose modular interfaces, while joint
generators permit sequence, backbone, and chemistry to adapt simultaneously.  The same trade-off
between controllability and identifiability reappears in reverse.

\subsection{Design as a choice of probabilistic factorization}

Let $x$ denote amino-acid sequence, $b$ backbone coordinates, $a$ all-atom coordinates including
side chains, $f$ a desired function, and $c$ other constraints such as a motif, target, symmetry, or
shape.  Contemporary design systems approximate different factorizations of
\begin{equation}
p(x,b,a,f\mid c).
\label{eq:design-joint}
\end{equation}
Three common strategies are
\begin{align}
\text{structure first:}\quad
&p(b\mid c)\,p(x\mid b,c)\,p(a_{\mathrm{side}}\mid x,b,c),
\label{eq:structure-first-design}\\
\text{sequence first:}\quad
&p(x\mid f,c)\,p(a\mid x,c),
\label{eq:sequence-first-design}\\
\text{joint generation:}\quad
&p(x,a\mid f,c).
\label{eq:joint-protein-design}
\end{align}
The first is modular and easy to audit, but errors accumulate across interfaces.  The third permits
side-chain and sequence decisions to reshape the backbone, but makes the source of a design signal
less identifiable.  In every case, function is only disentangled when $f$ is represented by an
explicit condition and tested against alternatives, not merely inferred from fold similarity.

\subsection{RFdiffusion and ProteinMPNN: a structure-first pipeline}

RFdiffusion fine-tunes RoseTTAFold as an SE(3)-aware denoiser over residue frames
\cite{watson2023}.  Translations are corrupted with Gaussian noise and rotations by Brownian motion
on $\mathrm{SO}(3)$.  Starting from noisy frames $B_T$, the reverse process samples
\begin{equation}
B_T\rightarrow B_{T-1}\rightarrow\cdots\rightarrow B_0
\sim p_{\theta}(b\mid c),
\label{eq:rfdiffusion}
\end{equation}
where $c$ may specify a fixed functional motif, target surface, fold, partial sequence, or symmetry.
Self-conditioning passes a previous denoised prediction into the next step.  The model normally
generates backbone geometry rather than amino-acid sequence or side-chain conformations.

ProteinMPNN then performs inverse folding on the sampled backbone \cite{dauparas2022}.  For a
decoding order $\pi$, its graph message-passing network parameterizes
\begin{equation}
p_{\phi}(x\mid b)
=\prod_{k=1}^{L}
p_{\phi}(x_{\pi_k}\mid b,x_{\pi_{<k}}).
\label{eq:proteinmpnn}
\end{equation}
Backbone atom distances and relative orientations are geometric inputs; residue identities are
categorical outputs.  Order-agnostic training, multichain graphs, tied positions, and positive or
negative weights on alternative backbones support symmetry and multistate design.

This pipeline has the clearest modular separation in current practice:
\begin{equation}
c\xrightarrow{\mathrm{RFdiffusion}}b
\xrightarrow{\mathrm{ProteinMPNN}}x
\xrightarrow{\mathrm{predictor}}\widehat b(x).
\label{eq:rfdiffusion-pipeline}
\end{equation}
Yet the separation is computational rather than biological.  A backbone motif is only a geometric
proxy for function; ProteinMPNN learns sequence compatibility from deposited structures, not
binding affinity or catalytic rate; and side chains cannot feed back while RFdiffusion chooses the
backbone.  Self-consistency with AlphaFold or ESMFold tests whether $x$ returns to $b$, but reuses
learned structural priors and is not an independent functional assay.

\subsection{Chroma: programmable posterior sampling}

Chroma also uses a structure-first factorization, but explicitly interprets controllable design as
Bayesian conditioning \cite{ingraham2023}.  A correlated diffusion prior reflects polymer chain
structure and radius-of-gyration statistics.  During sampling, external conditions modify the score
schematically as
\begin{equation}
\nabla_b\log p_t(b\mid c)
=\nabla_b\log p_t(b)
+\nabla_b\log p_t(c\mid b).
\label{eq:chroma-guidance}
\end{equation}
Conditions can encode symmetry, substructures, shape, distances, semantic classes, or language
annotations.  Random graph neural networks supply long-range communication with subquadratic
scaling, and a separate design network predicts sequence and side-chain conformations from the
sampled backbone.

Chroma separates the unconditional structural prior from user-provided constraint energies, which
makes combinations of controls transparent.  It does not guarantee that separately trained
conditioners are mutually compatible or calibrated.  Adding their score gradients is a
product-of-experts construction, not evidence that shape, function, and fold are independent.  Its
sequential backbone-then-sequence decoder also retains the same missing side-chain-to-backbone
feedback as Eq.~\ref{eq:structure-first-design}.

\subsection{MultiFlow: coupled discrete and continuous generation}

MultiFlow replaces the sequential factorization with a joint state for each residue
\cite{campbell2024}:
\begin{equation}
T_i=(r_i,R_i,x_i)
\in\R^3\times\mathrm{SO}(3)\times\mathcal A,
\label{eq:multiflow-state}
\end{equation}
where translations and rotations follow continuous flow matching, while amino-acid identities
follow a continuous-time Markov chain on a discrete alphabet.  A shared denoising network predicts
clean translations, rotations, and categorical probabilities.  By fixing one modality and flowing
the other, the same model can perform forward folding, inverse folding, or joint co-generation.

This is a mathematically clean track factorization: continuous geometry has a vector field and
discrete sequence has a rate matrix.  It removes the hard interface between a backbone generator
and a separately trained inverse-folding model.  However, the shared denoiser deliberately learns
cross-modal correlations, and the model state contains residue frames rather than explicit
side-chain chemistry.  Function, binding context, dynamics, and alternate conformations are not
separate variables in the basic objective.

\subsection{Protpardelle and Pallatom: two all-atom representations}

Protpardelle addresses the fact that amino-acid identity determines which side-chain atoms exist
\cite{chu2024}.  At every residue it maintains an ``atom73'' superposition containing backbone
atoms and the possible side-chain atoms of all twenty amino acids.  A structure denoiser updates
coordinates, while a non-autoregressive mini-ProteinMPNN predicts the current residue identity.
The selected identity collapses the superposition at each step:
\begin{equation}
\mathcal S_i^{(t)}
\xrightarrow{p_\phi(x_i\mid\mathcal S^{(t)})}
x_i^{(t)}
\xrightarrow{\mathrm{collapse}}
a_i^{(t)}
\xrightarrow{\mathrm{denoise}}
\mathcal S_i^{(t-1)}.
\label{eq:protpardelle-superposition}
\end{equation}
Side chains can therefore influence backbone updates during generation.  Nevertheless, sequence
has no independent diffusion process; Protpardelle is structure-primary, and its sequence estimate
acts as an iterative latent assignment.  The superposition solves changing dimensionality but is not
a physical ensemble.  Reported functional-group conditioning is promising, while conditional
fidelity remains weaker than backbone-only scaffolding.

Pallatom uses a different fixed-dimensional construction \cite{qu2024}.  Every residue is padded to
fourteen atom slots; atoms absent for the current amino acid are placed as virtual atoms at the
$C_\alpha$ coordinate.  An EDM-style Gaussian diffusion is learned on
\begin{equation}
a_t\in\R^{L\times14\times3},
\qquad
\mathcal L_{\mathrm{EDM}}
=\E_{a_0,a_t,\sigma}
\lambda(\sigma)\|D_\theta(a_t,\sigma)-a_0\|_2^2.
\label{eq:pallatom-loss}
\end{equation}
A dual-track network maintains globally communicating residue-token states and locally attending
atomic states.  Token information is broadcast to atoms, atomic information is aggregated back,
and denoised coordinates are recycled as pair and self-conditioning features.  A categorical head
decodes amino-acid identity from the generated atom geometry.

Pallatom deliberately learns $p(a)$ and relies on all-atom geometry to encode sequence and chemical
properties.  The atom14 padding removes changing dimensionality and reduces direct leakage from an
atom-existence mask, but it also embeds a representation convention into the coordinate
distribution.  Sequence, side-chain chemistry, packing, and backbone geometry are therefore not
disentangled; they are intentionally fused so that one coordinate denoiser can co-generate them.
This may improve compatibility, but a successful amino-acid readout does not show that catalytic
function, binding specificity, or physical energy has been separately learned.

\subsection{What the design models actually disentangle}

\begin{center}
\scriptsize
\begin{tabular}{>{\raggedright\arraybackslash}p{0.12\linewidth}
                >{\raggedright\arraybackslash}p{0.23\linewidth}
                >{\raggedright\arraybackslash}p{0.24\linewidth}
                >{\raggedright\arraybackslash}p{0.29\linewidth}}
\toprule
Model & Generated variables & Explicit separation & Remaining ambiguity \\
\midrule
RFdiffusion & residue-frame backbone & condition versus backbone; translation versus rotation & motif geometry versus biochemical function; no sequence or side chains \\
ProteinMPNN & sequence conditioned on a fixed backbone & geometric input versus categorical output & stability, expression, function, and negative states absent from ordinary training \\
Chroma & backbone, then sequence and side chains & structural prior versus external guidance; separate design decoder & conditioner calibration, product-of-experts conflicts, no sequence feedback during backbone diffusion \\
MultiFlow & residue frames and amino-acid sequence & continuous geometric flow versus discrete categorical flow & shared latent denoiser; side-chain chemistry and function absent \\
Protpardelle & backbone, side chains, and sequence estimate & atom superposition versus categorical collapse & sequence is not an independent generative process; superposition is algorithmic rather than physical \\
Pallatom & atom14 coordinates and decoded sequence & token-level versus atom-level neural tracks & geometry, residue identity, chemistry, and virtual-atom convention fused in $p(a)$ \\
\bottomrule
\end{tabular}
\end{center}

The design lesson is not that joint generation is always superior.  Modular factorization exposes
which component failed and permits specialized replacement; joint generation permits mutual
adaptation but makes attribution harder.  A controllable design system should therefore report
both interface-level controls and intervention-based evidence that changing one target property
does not silently change the others.

\section{Identification through interventions}

Architectural inspection establishes where information can flow, not what a hidden state means.
Meaning is identified only by comparing representations under perturbations that preserve declared
nuisances and remove a specific source of evidence.  The required intervention depends on whether
the source is present at inference, stored in pretrained parameters, supplied as a geometric track,
or imposed as a design condition.

\subsection{Sequence-axis nulls}

To identify MSA-specific modes, one should compare real alignments with null ensembles that
preserve selected marginals while destroying selected dependencies.  Useful controls include
independent-site shuffles, conservation-matched resampling, phylogeny-preserving simulations, and
row-identity perturbations.  For a representation $O(X)$, define
\begin{equation}
\Delta O
=O(X_{\mathrm{real}})
-\E_{\widetilde X\sim Q_{\mathrm{null}}}O(\widetilde X).
\label{eq:null-residual}
\end{equation}
Modes should be called coevolutionary only when they are enriched under an intervention that
removes pair dependence while retaining the relevant marginal and phylogenetic controls.

For AF1, the intervention is applied before feature generation.  For AF2, it can be followed
through $m_{si}$, outer-product-mean updates, and $z_{ij}$.  For AF3, the key measurement is how the
four-block MSA module changes $z_{ij}$ relative to the no-MSA or null-MSA condition.
For MSA Transformer, row and column interventions should be crossed: shuffle residues within
columns while preserving row phylogeny, then shuffle rows or simulate phylogeny while preserving
column marginals.  For RoseTTAFold, the same MSA nulls should be combined with template removal and
stopping feedback from the coordinate track.

ESM requires a different intervention because the training sequences are not present at inference.
With $\theta_{\mathrm{LM}}$ fixed, one can perturb the query sequence while matching composition,
motifs, or predicted secondary structure and measure changes in $h_i^{(\ell)}$ and
$A_{ij}^{\ell h}$.  To distinguish global evolutionary signal from database composition and
phylogeny, one must additionally compare models pretrained on reweighted data or controlled family
subsets.  For ESMFold, ablations should separately remove the attention contribution to
$z_{ij}^{(0)}$, alter the ESM layer mixture in $s_i^{(0)}$, and trace which modes are created or
suppressed by the folding trunk.

\subsection{Track and modality interventions}

For a multi-track model with output $O$ and information tracks $Y^1,\ldots,Y^K$, define the
conditional contribution of track $k$ relative to a matched null as
\begin{equation}
\Delta_k O
=O(Y^1,\ldots,Y^k,\ldots,Y^K)
-\E_{\widetilde Y^k\sim Q_k}
O(Y^1,\ldots,\widetilde Y^k,\ldots,Y^K).
\label{eq:track-intervention}
\end{equation}
In RoseTTAFold this separates MSA, template, pair, and coordinate feedback only if the null preserves
the scale and uncertainty of the removed state.  Zeroing a track can create an out-of-distribution
activation and overstate its importance.

In ESM3, one should measure the full conditional matrix
\begin{equation}
C_{ab}=I_\theta(Y^a;Y^b\mid Y^{\mathcal T\setminus\{a,b\}}),
\label{eq:cross-modal-information}
\end{equation}
or a predictive approximation to it, under modality dropout and matched prompts.  High
$C_{\mathrm{sequence},\mathrm{function}}$ shows cross-modal predictability, but not whether function
is mediated by structure, family identity, annotation pipelines, or training-set duplication.
Counterfactual prompts that hold structure fixed while changing function, or vice versa, are more
diagnostic than unconditional generation.

\subsection{Design-cycle validation and negative states}

Backbone design is commonly evaluated by generating sequences and folding them back:
\begin{equation}
b\xrightarrow{G_{\mathrm{seq}}}x
\xrightarrow{F_{\mathrm{fold}}}\widehat b,
\qquad
S_{\mathrm{cycle}}=\operatorname{sim}(b,\widehat b).
\label{eq:design-cycle}
\end{equation}
This is useful evidence of designability, but it is not an experimental likelihood and can be
inflated when the generator, inverse folder, and evaluator share training structures or learned
priors.  Strong evaluation should cross multiple sequence designers and folding models, search for
training-set neighbors, estimate uncertainty, and reserve biochemical function for independent
assays.

Functional design also requires negative states.  For a desired state $b^+$ and undesired folds,
complexes, or conformations $\{b_k^-\}$, a discriminative objective has the form
\begin{equation}
J(x)
=\log p(x\mid b^+)
-\sum_k\lambda_k\log p(x\mid b_k^-)
+\mu S_{\mathrm{func}}(x,b^+).
\label{eq:positive-negative-design}
\end{equation}
A model that optimizes only $b^+$ may create promiscuous binders, alternative oligomers, or
overstabilized inactive states.  Disentangled design should measure the cross-sensitivity matrix
between requested controls and realized properties rather than report only success on the prompted
property.

\subsection{Position-axis nuisance projection}

Let $M$ be a scalar position statistic decoded from a pair representation.  Let $Z$ contain known
position covariates such as conservation, entropy, gap rate, chain identity, relative position,
template coverage, and entity type.  A declared nuisance residual is
\begin{equation}
M_{\perp}
=Q_Z(M-M_0)Q_Z^*,
\qquad
Q_Z=I-Z(Z^{\mathsf T}WZ)^{-1}Z^{\mathsf T}W.
\label{eq:position-residual}
\end{equation}
This is preferable to deleting the first empirical eigenvector without specifying why it is
nuisance.  In AF3, entity-type stratification is mandatory because protein--protein,
protein--RNA, and protein--ligand pairs have different information sources.

\subsection{Causal and mechanistic limits}

An MSA-derived field is an evolutionary statistical object.  It can predict contacts and
epistasis without being a causal physical interaction.  A learned pair state is further removed
from direct interpretation because it combines evolutionary evidence, templates, sequence priors,
recycling, and geometric self-consistency.

Mechanistic claims require interventions or external endpoints: mutational epistasis, contact
maps, binding measurements, alternate conformations, or controlled removal of MSA and template
evidence.  Attention visualization and leading eigenmodes are diagnostic tools, not identification
theorems.

\section{Technical outlook and unresolved frontiers}

The next generation of protein models must move beyond predicting one deposited structure or
generating one self-consistent sequence--backbone pair.  A more complete target is
\begin{equation}
p\!\left(
x,\{a^{(k)},w_k\}_{k=1}^{K},
\Delta G,\kappa,f,e
\mid c\right),
\label{eq:future-protein-target}
\end{equation}
where $\{a^{(k)},w_k\}$ is a weighted conformational ensemble, $\Delta G$ denotes relative free
energies, $\kappa$ kinetic or transition quantities, $f$ molecular function, $e$ the biochemical or
cellular environment, and $c$ the design specification.  Current systems model only selected
marginals of Eq.~\ref{eq:future-protein-target}.

\subsection{Adaptive integration of family evidence and global priors}

Query MSAs preserve rare family- and subfamily-specific constraints, while language models provide
broad priors when homologues are absent.  Future architectures should retrieve evolutionary
evidence adaptively rather than choose a permanently MSA-based or MSA-free route.  A learned router
could retain row-level and subfamily modes only where they change the posterior relative to the
language-model prior, while compressing redundant regions early.  Typed sparse pair operators could
replace a uniformly dense $L\times L$ memory for long proteins and complexes.

The unresolved difficulty is calibration.  Database depth, taxonomic redundancy, family rarity,
and language-model confidence are measured on different scales.  Without an explicit likelihood
ratio or uncertainty model, adding retrieved homologues can overwrite a valid global prior, while a
large language model can smooth away a rare but decisive coupling.  Future hybrid models need to
report which prediction changed because of newly retrieved family evidence.

\subsection{From correlated modalities to causal biological variables}

ESM3-style modality tracks make sequence, structure, and function separately promptable, but the
tracks remain correlated through family identity and annotation pipelines.  Future multimodal
training should include counterfactual examples that hold structure fixed while varying function,
hold function fixed across unrelated scaffolds, and separate experimentally measured activities
from computational annotations.  Explicit latent variables for family, phylogeny, geometry,
chemical mechanism, and assay context would permit nuisance projection and intervention tests.

This will require new data rather than architecture alone.  Natural databases rarely contain the
factorial combinations needed for identification.  Negative biochemical results, weak binders,
inactive homologues, alternative substrates, and failed designs are especially valuable because
they break correlations that successful natural proteins preserve.

\subsection{Energy landscapes, ensembles, and kinetics}

Static coordinate prediction is not equivalent to thermodynamics.  Multiple diffusion samples are
also not automatically Boltzmann-weighted.  Future models need explicit representations of
conformational states and their relative populations, together with solvent, protonation,
electrostatics, metal coordination, membrane context, and temperature.  One possible decomposition
is
\begin{equation}
p(a\mid x,e)
=\sum_{k=1}^{K}w_k(x,e)\,p_k(a\mid x,e),
\qquad
w_k\propto\exp[-\beta G_k(x,e)].
\label{eq:ensemble-mixture}
\end{equation}
Learned generators could propose states, while differentiable physical models or experimental
observables calibrate $G_k$ and transition barriers.  This is essential for enzymes, receptors,
switches, transporters, and allosteric proteins, whose function depends on motion and population
shifts rather than one minimum-energy structure.

\subsection{Chemically typed, variable-topology all-atom generation}

Atom73 superposition and atom14 virtual padding convert variable side-chain topology into a fixed
tensor, but both expose artificial conventions to the network.  A longer-term solution is a
trans-dimensional generator over a typed chemical graph and continuous coordinates:
\begin{equation}
(V_t,E_t,A_t)
\longrightarrow
(V_{t-1},E_{t-1},A_{t-1}),
\label{eq:transdimensional-all-atom}
\end{equation}
where atoms and bonds can appear, disappear, or change type while coordinates evolve under an
equivariant flow.  Hard or differentiable constraints should preserve valence, chirality, bond
geometry, and excluded volume.  The same representation should extend beyond standard amino acids
to ligands, cofactors, covalent modifications, nucleic acids, ions, and waters.

The central unresolved issue is whether a learned coordinate distribution can recover chemical
energy and function without explicit atom types and electronic context.  All-atom self-consistency
demonstrates packing compatibility, but catalysis and specificity depend on charge states,
polarization, solvent rearrangement, and transition-state stabilization that are weakly represented
in structural databases.

\subsection{Specificity, negative design, and cellular developability}

Most design pipelines optimize a desired fold or interaction and only weakly represent competing
states.  Future systems should jointly model target binding, off-target binding, oligomerization,
aggregation, proteolysis, immunogenicity, expression, and localization.  Positive and negative
states should be first-class conditioning variables rather than filters applied after generation.

For functional proteins, the target is not simply a high structural confidence score.  Binding
design needs affinity and specificity; enzyme design needs reaction rate and substrate selectivity;
therapeutic design needs stability and behavior in a cellular or organismal environment.  These
outputs will require assay-specific uncertainty and cannot be reduced to a universal scalar fitness
without hiding trade-offs.

\subsection{Independent evaluation and experimental closed loops}

Current designability benchmarks often reuse ProteinMPNN, AlphaFold, ESMFold, or related training
distributions inside both generation and evaluation.  Future benchmarks should separate training
families temporally and structurally, use multiple independent folding and energy models, report
nearest-neighbor novelty, and reserve prospective experiments as the final endpoint.  Calibration
curves should relate model confidence to synthesis, expression, folding, binding, and activity
success rates.

The most consequential technical development may be a closed experimental loop:
\begin{equation}
\boxed{
\begin{gathered}
\text{retrieval and priors}
\rightarrow
\text{multistate all-atom design}
\rightarrow
\text{independent tests}\\
\longrightarrow
\text{multiplexed experiments}
\rightarrow
\text{uncertainty-aware model update}
\end{gathered}}
\label{eq:future-design-loop}
\end{equation}
Active learning should select experiments that distinguish competing mechanisms or reduce
uncertainty, rather than testing only the highest-scoring designs.  Data provenance, failed
experiments, assay conditions, and model versions must be retained so that improvements in scale
can be separated from improvements in biological identification.

\section{Conclusion}

Protein models are best understood as information-factorization systems.  They make two largely
independent choices: which evidence is acquired for the current query, and where that evidence is
compressed into sequence, pair, coordinate, or functional states.  Query MSAs and pretrained
language models are not competing definitions of evolution; they are family-specific and amortized
observations of evolutionary sequence distributions.  Pair tracks, attention maps, structure tokens,
and all-atom coordinates are subsequent computational representations of that evidence.

Protein design operates on the same joint distribution in reverse.  A modular pipeline factors
backbone generation, inverse folding, and side-chain construction; a joint generator allows those
variables to adapt together.  Modularity exposes failure boundaries but can create interface
inconsistency.  Joint generation reduces that inconsistency but intentionally mixes sequence,
geometry, and chemistry in a shared state.  Neither choice by itself establishes functional or
mechanistic understanding.

The main distinction throughout this review is therefore between \emph{computational
factorization} and \emph{biological disentanglement}.  Modern architectures increasingly separate
tensor axes, neural tracks, token vocabularies, output heads, and geometric symmetries.  Their
predictive power, however, often depends on combining phylogeny, conservation, structural
constraint, database composition, function labels, and learned physical priors.  A leading mode or
controllable prompt cannot be assigned one biological meaning from architecture alone.

Biological identification begins where architectural description ends.  It requires declared
categorical gauges, matched sequence and modality nulls, counterfactual track interventions,
negative design against competing states, calibrated uncertainty, and independent experimental
endpoints.  The most important future advance will not be a larger model that produces more
self-consistent proteins, but a model whose evidence, uncertainty, and causal claims remain
traceable from evolutionary data through all-atom design to experiment.

\scriptsize
\begin{multicols}{2}
\raggedcolumns
\begin{thebibliography}{99}
\setlength{\itemsep}{-0.2pt}
\setlength{\parskip}{0pt}

\bibitem{hoeffding1948}
W.~Hoeffding.
\newblock A class of statistics with asymptotically normal distribution.
\newblock \emph{Annals of Mathematical Statistics}, 19:293--325, 1948.

\bibitem{jones2012}
D.~T. Jones, D.~W. Buchan, D.~Cozzetto, and M.~Pontil.
\newblock PSICOV: precise structural contact prediction using sparse inverse covariance estimation
on large multiple sequence alignments.
\newblock \emph{Bioinformatics}, 28:184--190, 2012.

\bibitem{morcos2011}
F.~Morcos et al.
\newblock Direct-coupling analysis of residue coevolution captures native contacts across many
protein families.
\newblock \emph{PNAS}, 108:E1293--E1301, 2011.

\bibitem{ekeberg2013}
M.~Ekeberg et al.
\newblock Improved contact prediction in proteins: Using pseudolikelihoods to infer Potts models.
\newblock \emph{Physical Review E}, 87:012707, 2013.

\bibitem{cocco2018}
S.~Cocco, C.~Feinauer, M.~Figliuzzi, R.~Monasson, and M.~Weigt.
\newblock Inverse statistical physics of protein sequences: A key issues review.
\newblock \emph{Reports on Progress in Physics}, 81:032601, 2018.

\bibitem{dunn2008}
S.~D. Dunn, L.~M. Wahl, and G.~B. Gloor.
\newblock Mutual information without the influence of phylogeny or entropy dramatically improves
residue contact prediction.
\newblock \emph{Bioinformatics}, 24:333--340, 2008.

\bibitem{wang2024}
H.~Wang, S.~Feng, K.~Tsuboyama, S.~Liu, G.~J. Rocklin, and S.~Ovchinnikov.
\newblock Disentanglement of evolutionary constraints in statistical models of proteins.
\newblock \emph{PRX Life}, 2:023005, 2024.

\bibitem{senior2020}
A.~W. Senior et al.
\newblock Improved protein structure prediction using potentials from deep learning.
\newblock \emph{Nature}, 577:706--710, 2020.

\bibitem{jumper2021}
J.~Jumper et al.
\newblock Highly accurate protein structure prediction with AlphaFold.
\newblock \emph{Nature}, 596:583--589, 2021.

\bibitem{evans2022}
R.~Evans et al.
\newblock Protein complex prediction with AlphaFold-Multimer.
\newblock bioRxiv 2021.10.04.463034, revised 2022.

\bibitem{abramson2024}
J.~Abramson et al.
\newblock Accurate structure prediction of biomolecular interactions with AlphaFold 3.
\newblock \emph{Nature}, 630:493--500, 2024.

\bibitem{rao2021}
R.~Rao et al.
\newblock MSA Transformer.
\newblock \emph{ICML}, 2021.

\bibitem{baek2021}
M.~Baek et al.
\newblock Accurate prediction of protein structures and interactions using a three-track neural
network.
\newblock \emph{Science}, 373:871--876, 2021.

\bibitem{rives2021}
A.~Rives et al.
\newblock Biological structure and function emerge from scaling unsupervised learning to 250
million protein sequences.
\newblock \emph{PNAS}, 118:e2016239118, 2021.

\bibitem{rao2020}
R.~Rao et al.
\newblock Transformer protein language models are unsupervised structure learners.
\newblock bioRxiv 2020.12.15.422761, 2020.

\bibitem{meier2021}
J.~Meier et al.
\newblock Language models enable zero-shot prediction of the effects of mutations on protein
function.
\newblock \emph{NeurIPS}, 2021.

\bibitem{lin2023}
Z.~Lin et al.
\newblock Evolutionary-scale prediction of atomic-level protein structure with a language model.
\newblock \emph{Science}, 379:1123--1130, 2023.

\bibitem{hayes2024}
T.~Hayes et al.
\newblock Simulating 500 million years of evolution with a language model.
\newblock \emph{Science}, 387:850--858, 2025.

\bibitem{watson2023}
J.~L. Watson et al.
\newblock De novo design of protein structure and function with RFdiffusion.
\newblock \emph{Nature}, 620:1089--1100, 2023.

\bibitem{dauparas2022}
J.~Dauparas et al.
\newblock Robust deep learning-based protein sequence design using ProteinMPNN.
\newblock \emph{Science}, 378:49--56, 2022.

\bibitem{ingraham2023}
J.~Ingraham et al.
\newblock Illuminating protein space with a programmable generative model.
\newblock \emph{Nature}, 623:1070--1078, 2023.

\bibitem{campbell2024}
A.~Campbell, J.~Yim, R.~Barzilay, T.~Rainforth, and T.~Jaakkola.
\newblock Generative flows on discrete state-spaces: Enabling multimodal flows with applications
to protein co-design.
\newblock \emph{ICML}, 2024.

\bibitem{chu2024}
A.~E. Chu et al.
\newblock An all-atom protein generative model.
\newblock \emph{PNAS}, 121:e2311500121, 2024.

\bibitem{qu2024}
W.~Qu, J.~Guan, R.~Ma, K.~Zhai, W.~Wu, and H.~Wang.
\newblock P(all-atom) is unlocking new path for protein design.
\newblock \emph{Proceedings of the 42nd International Conference on Machine Learning},
PMLR 267:50786--50816, 2025.

\end{thebibliography}
\end{multicols}

\end{document}
